\documentclass[10pt,a4paper]{amsart}
\usepackage{amsmath,amssymb,amsthm,graphicx}
\usepackage{verbatim}
\newenvironment{reference}{\comment}{\endcomment}
%\newenvironment{reference}{}{}
\newenvironment{slogan}{\comment}{\endcomment}
\newenvironment{history}{\comment}{\endcomment}
% For commutative diagrams we use Xy-pic
\usepackage[all]{xy}

% We use 2cell for 2-commutative diagrams.
\xyoption{2cell}
% Theorem environments.
%
\theoremstyle{plain}
\newtheorem{theorem}[subsection]{Theorem}
\newtheorem{proposition}[subsection]{Proposition}
\newtheorem{lemma}[subsection]{Lemma}

\theoremstyle{definition}
\newtheorem{definition}[subsection]{Definition}
\newtheorem{example}[subsection]{Example}
\newtheorem{exercise}[subsection]{Exercise}
\newtheorem{situation}[subsection]{Situation}

\theoremstyle{remark}
\newtheorem{remark}[subsection]{Remark}
\newtheorem{remarks}[subsection]{Remarks}

\numberwithin{equation}{subsection}

% Macros
%
\def\lim{\mathop{\mathrm{lim}}\nolimits}
\def\colim{\mathop{\mathrm{colim}}\nolimits}
\def\Spec{\mathop{\mathrm{Spec}}}
\def\Hom{\mathop{\mathrm{Hom}}\nolimits}
\def\Ext{\mathop{\mathrm{Ext}}\nolimits}
\def\SheafHom{\mathop{\mathcal{H}\!\mathit{om}}\nolimits}
\def\SheafExt{\mathop{\mathcal{E}\!\mathit{xt}}\nolimits}
\def\Sch{\mathit{Sch}}
\def\Mor{\mathop{\mathrm{Mor}}\nolimits}
\def\Ob{\mathop{\mathrm{Ob}}\nolimits}
\def\Sh{\mathop{\mathit{Sh}}\nolimits}
\def\NL{\mathop{N\!L}\nolimits}
\def\CH{\mathop{\mathrm{CH}}\nolimits}
\def\proetale{{pro\text{-}\acute{e}tale}}
\def\etale{{\acute{e}tale}}
\def\QCoh{\mathit{QCoh}}
\def\Ker{\mathop{\mathrm{Ker}}}
\def\Im{\mathop{\mathrm{Im}}}
\def\Coker{\mathop{\mathrm{Coker}}}
\def\Coim{\mathop{\mathrm{Coim}}}

% Boxtimes
%
\DeclareMathSymbol{\boxtimes}{\mathbin}{AMSa}{"02}

%
% Macros for moduli stacks/spaces
%
\def\QCohstack{\mathcal{QC}\!\mathit{oh}}
\def\Cohstack{\mathcal{C}\!\mathit{oh}}
\def\Spacesstack{\mathcal{S}\!\mathit{paces}}
\def\Quotfunctor{\mathrm{Quot}}
\def\Hilbfunctor{\mathrm{Hilb}}
\def\Curvesstack{\mathcal{C}\!\mathit{urves}}
\def\Polarizedstack{\mathcal{P}\!\mathit{olarized}}
\def\Complexesstack{\mathcal{C}\!\mathit{omplexes}}
% \Pic is the operator that assigns to X its picard group, usage \Pic(X)
% \Picardstack_{X/B} denotes the Picard stack of X over B
% \Picardfunctor_{X/B} denotes the Picard functor of X over B
\def\Pic{\mathop{\mathrm{Pic}}\nolimits}
\def\Picardstack{\mathcal{P}\!\mathit{ic}}
\def\Picardfunctor{\mathrm{Pic}}
\def\Deformationcategory{\mathcal{D}\!\mathit{ef}}
\usepackage{cleveref}
\crefname{originalreference}{original reference}{original references}
\Crefname{originalreference}{Original reference}{Original references}
\makeatletter
\newlabel{algebra-remark-Tor-ring-mod-ideal}{{algebra-remark-Tor-ring-mod-ideal}{source}{}{sourcepointer-1}{}}
\newlabel{algebra-remark-Tor-ring-mod-ideal@cref}{{[originalreference][0][]algebra-remark-Tor-ring-mod-ideal}{[1][0][]source}{}{}{}}
\makeatother
\crefname{originalreference}{original reference}{original references}
\Crefname{originalreference}{Original reference}{Original references}
\makeatletter
\newlabel{more-algebra-lemma-tor-strictly-pro-zero}{{more-algebra-lemma-tor-strictly-pro-zero}{source}{}{sourcepointer-2}{}}
\newlabel{more-algebra-lemma-tor-strictly-pro-zero@cref}{{[originalreference][0][]more-algebra-lemma-tor-strictly-pro-zero}{[1][0][]source}{}{}{}}
\makeatother
\title{Maps between Tor modules of a module killed by a power of the ideal vanish after shifting the index far enough}
\author{The Stacks Project Authors}
\date{stacks-project:local-cohomology.tex:local-snapshot-20261012; GNU FDL 1.2 or later}
\begin{document}
\maketitle
\noindent\textit{Editorial note.} All delivered mathematical text is the author's own, copied from the pinned Stacks Project chapter source sources/stacks-local-cohomology/pinned-original-source.tex; no statement or proof step was rewritten, reordered or removed.\par\medskip
\noindent\textit{Editorial note.} Preamble provenance: the Stacks Project keeps one single shared preamble for every chapter; this card takes it directly from the staged sources/stacks-local-cohomology/preamble.tex as three verbatim spans (lines 8-12, 14-18 and 154-217), so no derived file was written and nothing was added or redefined. The wrapper class and the base packages are supplied by the assembler.\par\medskip
\noindent\textit{Editorial note.} Bibliography provenance: neither the counted statement nor its proof contains a citation (this lane rescans the whole span before writing the card), so the card carries no bibliography and no entry was reconstructed.\par\medskip
\noindent\textit{Editorial note.} Selection provenance: this card was opened under the file-based control contract (CONTROL.md), with the mechanical pre-card gate run on its statement and proof interval before the card was written.\par\medskip
\noindent\textit{Allowed standard prerequisites (editorial).}\begin{itemize}
\item The identification of the first Tor module of the residue ring with a submodule of the tensor product of the ideal power with the module. Invoked inside the counted proof at source lines 4521-4527 and attributed to a remark of the cited chapter (a cross-chapter pointer). Provides the base case of the induction.
\item The vanishing of Tor modules for a module annihilated by a power of the ideal in strictly positive degrees. Used inside the counted proof in the induction step, attributed to a lemma of the cited chapter (a cross-chapter pointer). Supplies the input that makes the induction step work.
\end{itemize}\medskip
\ifdefined\hypertarget\hypertarget{sourcepointer-1}{}\fi
\noindent\textit{Source reference.} Original reference algebra-remark-Tor-ring-mod-ideal: sources/stacks-local-cohomology/pinned-original-source.tex, source lines 4525--4525. Navigation pointer only: the counted text refers to the Stacks item algebra-remark-Tor-ring-mod-ideal, which belongs to another chapter of the project and is not part of this card. The pointer carries the target's own label as its number because this chapter's staged snapshot does not contain the other chapter's tag table; its content is not reproduced here.\par\medskip
\ifdefined\hypertarget\hypertarget{sourcepointer-2}{}\fi
\noindent\textit{Source reference.} Original reference more-algebra-lemma-tor-strictly-pro-zero: sources/stacks-local-cohomology/pinned-original-source.tex, source lines 4492--4492. Navigation pointer only: the counted text refers to the Stacks item more-algebra-lemma-tor-strictly-pro-zero, which belongs to another chapter of the project and is not part of this card. The pointer carries the target's own label as its number because this chapter's staged snapshot does not contain the other chapter's tag table; its content is not reproduced here.\par\medskip
\begin{lemma}
\label{lemma-maps-zero-fixed-torsion}
Let $I$ be an ideal of a Noetherian ring $A$. For every $m \geq 0$
and $i > 0$ there exist a $c = c(A, I, m, i) \geq 0$ such that
for every $A$-module $M$ annihilated by $I^m$ the map
$$
\text{Tor}^A_i(M, A/I^n) \to \text{Tor}^A_i(M, A/I^{n - c})
$$
is zero for all $n \geq c$.
\end{lemma}

\begin{proof}
By induction on $i$. Base case $i = 1$. The short exact sequence
$0 \to I^n \to A \to A/I^n \to 0$ determines an injection
$\text{Tor}_1^A(M, A/I^n) \subset I^n \otimes_A M$, see
Algebra, Remark \ref{algebra-remark-Tor-ring-mod-ideal}.
As $M$ is annihilated by $I^m$ we see that the map
$I^n \otimes_A M \to I^{n - m} \otimes_A M$ is
zero for $n \geq m$. Hence the result holds with $c = m$.

\medskip\noindent
Induction step. Let $i > 1$ and assume $c$ works for $i - 1$.
By More on Algebra, Lemma \ref{more-algebra-lemma-tor-strictly-pro-zero}
applied to $M = A/I^m$ we can choose $c' \geq 0$ such that
$\text{Tor}_i(A/I^m, A/I^n) \to \text{Tor}_i(A/I^m, A/I^{n - c'})$
is zero for $n \geq c'$. Let $M$ be annihilated by $I^m$. Choose a short
exact sequence
$$
0 \to S \to \bigoplus\nolimits_{i \in I} A/I^m \to M \to 0
$$
The corresponding long exact sequence of tors gives an exact sequence
$$
\text{Tor}_i^A(\bigoplus\nolimits_{i \in I} A/I^m, A/I^n) \to
\text{Tor}_i^A(M, A/I^n) \to
\text{Tor}_{i - 1}^A(S, A/I^n)
$$
for all integers $n \geq 0$. If $n \geq c + c'$, then the map
$\text{Tor}_{i - 1}^A(S, A/I^n) \to \text{Tor}_{i - 1}^A(S, A/I^{n - c})$
is zero and the map $\text{Tor}_i^A(A/I^m, A/I^{n - c}) \to 
\text{Tor}_i^A(A/I^m, A/I^{n - c - c'})$ is zero. Combined with the
short exact sequences this implies the result holds for $i$ with
constant $c + c'$.
\end{proof}

\end{document}
